\begin{frame}{¿Qué son las apropiaciones?}
\textbf{Apropiaciones:} son los \textbf{montos máximos autorizados por la Ley de Presupuesto} para ser comprometidos y ejecutados por una entidad pública.

\vspace{1em}
\begin{itemize}
  \item Representan el límite legal de gasto por cada rubro presupuestal.
  \item No implican ejecución automática ni gasto efectivo.
  \item Solo el Congreso puede modificarlas.
\end{itemize}

\vspace{1em}
\textit{Son el punto de partida del ciclo presupuestal.}
\end{frame}

\begin{frame}{¿Qué son los compromisos?}
\textbf{Compromisos:} son actos administrativos que generan una obligación futura de pago por parte del Estado.

\vspace{1em}
\begin{itemize}
  \item Se generan al contratar o adquirir bienes y servicios.
  \item Disminuyen el saldo disponible de la apropiación.
  \item Aún no son pagos, pero sí obligaciones asumidas formalmente.
\end{itemize}

\vspace{1em}
\textit{Ejemplo: firma de contrato de obra pública.}
\end{frame}

\begin{frame}{¿Qué son las obligaciones?}
\textbf{Obligaciones:} son los montos reconocidos como deuda cierta, exigible y cuantificada a favor de un tercero.

\vspace{1em}
\begin{itemize}
  \item Surgen cuando el bien o servicio ha sido entregado o prestado.
  \item Se sustentan con factura o documento equivalente.
  \item Deben pagarse dentro de los plazos legales.
\end{itemize}

\vspace{1em}
\textit{Ejemplo: recepción de una obra terminada.}
\end{frame}

\begin{frame}{¿Qué son los pagos?}
\textbf{Pagos:} son las erogaciones efectivas de recursos públicos para cumplir con obligaciones previamente reconocidas.

\vspace{1em}
\begin{itemize}
  \item Representan el último paso del ciclo presupuestal.
  \item Requieren disponibilidad de caja (tesorería).
  \item Pueden realizarse en el mismo año o en vigencias futuras.
\end{itemize}

\vspace{1em}
\textit{Ejemplo: giro de fondos al contratista.}
\end{frame}

\begin{frame}{Ciclo de ejecución presupuestal en Colombia}
\begin{center}
\begin{tikzpicture}[
  node distance=0.7cm,
  every node/.style={draw, rectangle, rounded corners, minimum width=3.8cm, minimum height=1cm, align=center},
  arrow/.style={->, thick}
]

\node (apropiacion) {Apropiación\\(autorización legal de gasto)};
\node (compromiso) [below=of apropiacion] {Compromiso\\(acto administrativo que genera obligación)};
\node (obligacion) [below=of compromiso] {Obligación\\(reconocimiento de deuda exigible)};
\node (pago) [below=of obligacion] {Pago\\(salida efectiva de recursos)};

\draw[arrow] (apropiacion) -- (compromiso);
\draw[arrow] (compromiso) -- (obligacion);
\draw[arrow] (obligacion) -- (pago);

\end{tikzpicture}
\end{center}
\end{frame}


\begin{frame}{Ciclo de ejecución presupuestal en Colombia}
\vspace{1em}
\begin{itemize}
  \item El ciclo puede abarcar más de una vigencia fiscal.
  \item Solo los pagos generan una salida de caja efectiva.
  \item En general, solo se considere como gasto cuando los recursos son \textbf{obligados}
\end{itemize}
\end{frame}

\begin{frame}{¿Qué es la ejecución presupuestal?}
\textbf{Ejecución presupuestal:} es el grado en que las apropiaciones autorizadas en el presupuesto se convierten en obligaciones durante el año fiscal.

\vspace{1em}
\begin{itemize}
  \item Mide cuánto del presupuesto aprobado va a ser irreversiblemente gastado.
  \item Se expresa usualmente como el porcentaje de obligaciones frente a apropiaciones.
\end{itemize}

\vspace{1em}
\textbf{Fórmula común:}
\begin{equation*}
\text{Ejecución (\%)} = \frac{\text{Obligaciones}}{\text{Apropiaciones}} \times 100
\end{equation*}

\vspace{1em}
\textit{Una alta ejecución no siempre implica buena calidad del gasto.}
\end{frame}

\begin{frame}{¿Qué pasa si una apropiación no se compromete en la vigencia?}
\begin{itemize}
  \item Las \textbf{apropiaciones} son válidas solo durante la vigencia fiscal (normalmente un año).
  \item Si al cierre de la vigencia una apropiación \textbf{no se compromete}, esa parte del presupuesto \textbf{se pierde} (pérdida de apropiación).
  \item En consecuencia, los recursos \textbf{no comprometidos se liberan o se cancelan} y se pierde la posibilidad de uso.
\end{itemize}
\end{frame}

\begin{frame}{¿Qué pasa si un compromiso no se obliga en la vigencia?}
\begin{itemize}
  \item Si existe un \textbf{compromiso registrado pero no obligado}, puede convertirse en \textbf{reserva presupuestal}.
  \item La reserva permite hacer el pago en el año siguiente, siempre que se haya comprometido oportunamente.
\end{itemize}
\end{frame}

\begin{frame}
\textbf{En resumen:}
\vspace{1em}
\begin{itemize}
  \item \textbf{Apropiación no comprometido} → se pierde.
  \item \textbf{Compromiso no obligado} → se convierte en \textbf{reserva presupuestal}.
\end{itemize}
\end{frame}



\begin{frame}
\vspace{1em}
\textbf{Implicaciones:}
\begin{itemize}
  \item Las entidades buscan alcanzar a al menos comprometer los recursos para no perderlos.
  \item Desde el punto de vista del cierre fiscal de la vigencia, una menor ejecución (las apropiaciones no alcanzan a obligarse) representa una mejora en el balance fiscal (aunque las reservas representan mayor presión fiscal para la vigencia siguiente).
  \item El retraso en los pagos facilita el manejo de caja.   
\end{itemize}

\end{frame}

\begin{frame}{Organigrama simplificado del Ministerio de Hacienda}
\begin{center}
\begin{tikzpicture}[
  every node/.style={rectangle, draw, align=center, rounded corners, font=\scriptsize, minimum width=3.2cm, minimum height=0.8cm},
  node distance=0.6cm
]

% Nivel 1
\node (ministro) {Ministro de Hacienda};

% Nivel 2
\node[below left=1.2cm and 1.2cm of ministro] (vicgen) {Viceministerio\\General};
\node[below right=1.2cm and 1.2cm of ministro] (victec) {Viceministerio\\Técnico};

% Nivel 3 - General
\node[below=of vicgen] (dgp) {DGPPN\\Presupuesto};
\node[below=of dgp] (dgcptn) {DGCPTN\\Crédito y Tesoro};
\node[below=of dgcptn] (dgaf) {DGAF\\Apoyo Fiscal};

% Nivel 3 - Técnico
\node[below=of victec] (polmacro) {DGPM \\ Política \\ Macroeconómica};
\node[below=of polmacro] (dgrees) {DGRESS\\Seguridad Social};

% Conexiones
\draw[->] (ministro) -- (vicgen);
\draw[->] (ministro) -- (victec);

\draw[->] (vicgen) -- (dgp);
\draw[->] (dgp) -- (dgcptn);
\draw[->] (dgcptn) -- (dgaf);

\draw[->] (victec) -- (polmacro);
\draw[->] (polmacro) -- (dgrees);

\end{tikzpicture}
\end{center}
\end{frame}



